% SPDX-License-Identifier: 0BSD
\section{\texorpdfstring{Shared-Current DAG and Numerical Execution}{Shared-Current DAG and Numerical Execution}}
\label{sec:zgg-dag}

The current recursion becomes a directed acyclic graph (DAG) by merging
physically identical intermediate currents.  This section describes the
identity, staging, and timing conventions without depending on a particular
artifact format or software interface.

\subsection{Current identity}

Two contributions share a current only when every physics-relevant entry in
\begin{equation}
 \Omega(J)=\bigl(
 t,\ I,\ H,\ \chi,\ s,\ F,\ Q,\ C,\ O,\ a,\ \pi_I
 \bigr)
 \label{eq:zgg-current-key}
\end{equation}
is equal.  These entries specify the particle species \(t\), external subset
\(I\), helicity ancestry \(H\), chirality and spin state \((\chi,s)\), flavour
and additive quantum-number flow \((F,Q)\), colour state \(C\), perturbative
order \(O\), any auxiliary-current type \(a\), and the colour-sensitive
ordering \(\pi_I\).  The subset fixes the current momentum
\(q_I=\sum_{i\in I}q_i\); the ordered labels prevent distinct colour
traversals from being merged.

Exact equality permits reuse across all amplitude contributions and physical
flows.  A second level of reuse applies when two target currents require the
same kernel value up to a known sign or scalar factor.  It is useful to keep
the three concepts distinct:
\begin{equation}
 \underbrace{\text{same current identity}}_{\text{one stored node}},\qquad
 \underbrace{\text{same kernel evaluation}}_{\text{one numerical value}},\qquad
 \underbrace{\text{same physical flow}}_{\text{one resolved component}}.
 \label{eq:zgg-three-identities}
\end{equation}

\subsection{Topological stages and storage}

The graph is evaluated in increasing source-subset size.  An interaction
forming a current from \(k\) external labels belongs to stage \(k\), so every
input is available before its consumer.  For the process of
Sec.~\ref{sec:zgg-example}, the stage sequence is schematically
\begin{equation}
 \boxed{\text{sources}}
 \longrightarrow
 \boxed{|I|=2}
 \longrightarrow
 \boxed{|I|=3}
 \longrightarrow
 \boxed{|I|=4}
 \longrightarrow
 \boxed{\text{amplitude roots}} .
 \label{eq:zgg-stage-chain}
\end{equation}
At each stage, every admissible partition contributing to the same current is
summed coherently before the current is propagated. The amplitude roots then
close the largest retained currents against the remaining external states.
This ordering makes every dependency explicit without assigning significance
to storage locations or evaluator-specific identifiers.

A current may be needed as a source value, before propagation, or after
propagation.  Keeping these storage variants explicit prevents a propagator
from being applied twice and lets the final closure request the correct form.
Conceptually, the numerical state is
\begin{equation}
  x=\bigl[x_{\rm current}\mid x_{\rm momentum}\mid x_{\rm model}\bigr].
  \label{eq:zgg-global-state}
\end{equation}
Each stage consumes only the entries it needs:
\begin{equation}
  x_k=L_k x,\qquad y_k=E_k(x_k),\qquad
  x[S_k]\mathrel{\leftarrow}y_k ,
  \label{eq:zgg-stage-map}
\end{equation}
where \(L_k\) gathers the stage-local inputs and \(S_k\) identifies its output
slots.  This stage-local layout reduces data movement while leaving the
physics graph unchanged.

\subsection{Compiled and eager evaluation}

In compiled mode, a stage is optimized as a multi-output expression.  For
target current \(j\),
\begin{equation}
  y_j=\mathcal P_j\!\left[
       \sum_{e:\,\operatorname{target}(e)=j}
       f_e\,K_{g(e)}(x_k)
     \right],
  \label{eq:zgg-stage-expression}
\end{equation}
where \(K_{g(e)}\) is a reusable kernel value, \(f_e\) is the exact attachment
factor, and \(\mathcal P_j\) is the propagator or identity.  Symbolic
optimization may therefore share algebra across several current components
within the stage.

Eager mode evaluates the same graph with prepared model kernels and explicit
gather, accumulation, finalization, and closure steps.  Recurrence mode builds
a compact current schedule directly.  The three modes preserve the same stage
dependencies, coherent root sums, colour contraction, symmetry weights, and
normalization; only the placement of symbolic compilation and numerical
scheduling differs.

For leading-colour resolved evaluation, the public tensor is
\begin{equation}
  R_{n,h,c}=|{\cal M}_{h,c}(p_n)|^2,
  \qquad
  \operatorname{shape}(R)=
  (N_{\rm point},N_{\rm helicity},N_{\rm flow}) .
  \label{eq:zgg-resolved-shape}
\end{equation}
Structural zeros remain explicit.  Summing the resolved tensor reproduces the
directly evaluated total within the validation tolerance.

\subsection{Timing boundary}

The report separates the complete native evaluation from its numerical-kernel
submetric.  The complete time begins after caller-language array conversion and
includes source construction, momentum and parameter preparation, stage input
gathers, kernel execution, output assignment, closures, and final reduction.
The narrower execution time contains the mode-specific numerical kernels:
\begin{equation}
 t_{\rm exec}=\sum_{k}t\bigl(E_k\bigr)+t\bigl(E_{\rm root}\bigr),
 \qquad
 t_{\rm wall}=t_{\rm exec}+t_{\rm source}+t_{\rm momentum/model}
 +t_{\rm gather/scatter}+t_{\rm reduce}+t_{\rm batch}.
 \label{eq:zgg-timing-boundary}
\end{equation}
This distinction exposes orchestration overhead while keeping the complete
wall-time observable comparable across execution modes.
